% --- Archon physics formalization source begin ---
% archon:physics
% archon:covers PhyXMiniProblems/problem_phyx_mini_0086.lean
% archon:source-report reports/phyx_mini/problem_phyx_mini_0086.source.json

\chapter{Physics problem phyx\_mini\_0086}
\label{ch:PhyXMiniProblems_problem_phyx_mini_0086}

\paragraph{Problem source.}
\#\# Physical scenario

Figure shows an optical fiber in which a central plastic core of index of refraction \textbackslash{}( n\_1 = 1.58 \textbackslash{}) is surrounded by a plastic sheath of index of refraction \textbackslash{}( n\_2 = 1.53 \textbackslash{}). Light can travel along different paths within the central core, leading to different travel times through the fiber. This causes an initially short pulse of light to spread as it travels along the fiber, resulting in information loss.

\#\# Question

Consider light that travels directly along the central axis of the fiber and light that is repeatedly reflected at the critical angle along the core-sheath interface, reflecting from side to side as it travels down the central core. If the fiber length is \textbackslash{}( 300\textbackslash{},\textbackslash{}text\{m\} \textbackslash{}), what is the difference in the travel times along these two routes?

\#\# Answer choices

- A: A: 5.16ns

- B: B: 56.1ns

- C: C: 51.6ns

- D: D: 46.5ns

\#\# Figure caption (auxiliary; use the image as primary evidence)

The image depicts a graphical representation of a light ray traveling through an optical fiber. The fiber consists of a core with refractive index \textbackslash{}( n\_1 \textbackslash{}) (illustrated in light blue) and a cladding with a lower refractive index \textbackslash{}( n\_2 \textbackslash{}) (illustrated in light orange). The light path is indicated by a red line with arrows, showing the light entering at an angle, reflecting off the core-cladding boundary, and continuing through the core. The dashed black lines represent the axis of the cylindrical fiber structure.

\#\# Dataset metadata

Geometrical Optics; Spatial Relation Reasoning, Multi-Formula Reasoning

\paragraph{Recorded answer/context.}
C: C: 51.6ns

\paragraph{Figure/image path.}
/root/proposal\_for\_physic/science-mango/phyx\_mini\_run/phyx\_data/test\_image/86.png

\paragraph{Formalization target.}
create a compiling Lean file with sorry bodies at `PhyXMiniProblems/problem\_phyx\_mini\_0086.lean`. The Lean declarations must preserve the physical quantities, dimensions or dimensional roles, figure labels, governing-law hypotheses, and final relation expressed by this problem.
Use Mathlib/Physlib names found through LeanExplore where available. If a physics API is missing, introduce faithful local abstractions rather than scalar placeholder aliases.

\begin{theorem}[Physics formalization target]
\label{thm:physics:phyx_mini_0086:target}
\lean{PhyXMiniProblems.ProblemPhyXMini0086.optical_fiber_travel_time_difference}
\uses{def:physics:phyx-mini-0086:phyxminiproblems-problemphyxmini0086-opticalfiberpulsesetup, def:physics:phyx-mini-0086:phyxminiproblems-problemphyxmini0086-lengthinmeters, def:physics:phyx-mini-0086:phyxminiproblems-problemphyxmini0086-speedinmeterspersecond, def:physics:phyx-mini-0086:phyxminiproblems-problemphyxmini0086-traveltimedifferenceseconds, def:physics:phyx-mini-0086:phyxminiproblems-problemphyxmini0086-traveltimedifferencenanoseconds, def:physics:phyx-mini-0086:phyxminiproblems-problemphyxmini0086-matchesopticalfiberfigure, def:physics:phyx-mini-0086:phyxminiproblems-problemphyxmini0086-satisfiesopticalfiberlaws, def:physics:phyx-mini-0086:phyxminiproblems-problemphyxmini0086-isclosestanswerchoice}
The assigned autoformalize agent should translate this physics problem into Lean declarations in the covered file, with theorem and lemma proof bodies written as `by sorry`.
\end{theorem}
\begin{proof}
This is an autoformalization task, not a proof task. Produce statements that can later be proved by the physics prover without weakening the physical model.
\end{proof}

% --- Archon declaration topology begin ---
\section{Declaration topology}
\begin{definition}[Dim Length]
  \label{def:physics:phyx-mini-0086:phyxminiproblems-problemphyxmini0086-dimlength}
  \lean{PhyXMiniProblems.ProblemPhyXMini0086.DimLength}
  A physical length, independent of any particular choice of units
\end{definition}
\begin{proof}
  This is the declared model definition of Dim Length.
\end{proof}
\begin{definition}[Dim Time]
  \label{def:physics:phyx-mini-0086:phyxminiproblems-problemphyxmini0086-dimtime}
  \lean{PhyXMiniProblems.ProblemPhyXMini0086.DimTime}
  A physical time interval, independent of any particular choice of units
\end{definition}
\begin{proof}
  This is the declared model definition of Dim Time.
\end{proof}
\begin{definition}[Dim Speed Real]
  \label{def:physics:phyx-mini-0086:phyxminiproblems-problemphyxmini0086-dimspeedreal}
  \lean{PhyXMiniProblems.ProblemPhyXMini0086.DimSpeedReal}
  A physical speed, independent of any particular choice of units
\end{definition}
\begin{proof}
  This is the declared model definition of Dim Speed Real.
\end{proof}
\begin{definition}[nanosecond Unit Choices]
  \label{def:physics:phyx-mini-0086:phyxminiproblems-problemphyxmini0086-nanosecondunitchoices}
  \lean{PhyXMiniProblems.ProblemPhyXMini0086.nanosecondUnitChoices}
  Unit choices in which time readouts are in nanoseconds and other units are SI
\end{definition}
\begin{proof}
  This is the declared model definition of nanosecond Unit Choices.
\end{proof}
\begin{definition}[Optical Fiber Region]
  \label{def:physics:phyx-mini-0086:phyxminiproblems-problemphyxmini0086-opticalfiberregion}
  \lean{PhyXMiniProblems.ProblemPhyXMini0086.OpticalFiberRegion}
  The two homogeneous plastic regions identified by n₁ and n₂ in the figure
\end{definition}
\begin{proof}
  This is the declared model definition of Optical Fiber Region.
\end{proof}
\begin{definition}[Fiber Route]
  \label{def:physics:phyx-mini-0086:phyxminiproblems-problemphyxmini0086-fiberroute}
  \lean{PhyXMiniProblems.ProblemPhyXMini0086.FiberRoute}
  The two routes compared in the question
\end{definition}
\begin{proof}
  This is the declared model definition of Fiber Route.
\end{proof}
\begin{definition}[Optical Fiber Pulse Setup]
  \label{def:physics:phyx-mini-0086:phyxminiproblems-problemphyxmini0086-opticalfiberpulsesetup}
  \lean{PhyXMiniProblems.ProblemPhyXMini0086.OpticalFiberPulseSetup}
  \uses{def:physics:phyx-mini-0086:phyxminiproblems-problemphyxmini0086-dimlength, def:physics:phyx-mini-0086:phyxminiproblems-problemphyxmini0086-dimtime, def:physics:phyx-mini-0086:phyxminiproblems-problemphyxmini0086-dimspeedreal, def:physics:phyx-mini-0086:phyxminiproblems-problemphyxmini0086-opticalfiberregion, def:physics:phyx-mini-0086:phyxminiproblems-problemphyxmini0086-fiberroute}
  The quantities belonging to the fiber and the two depicted light routes. criticalIncidenceAngleRadians is measured from the normal to the core--sheath interface. criticalRayAngleToAxisRadians is measured from the dashed central axis. No numerical value for the requested travel-time difference is stored in this structure
\end{definition}
\begin{proof}
  This is the declared model definition of Optical Fiber Pulse Setup.
\end{proof}
\begin{definition}[length In Meters]
  \label{def:physics:phyx-mini-0086:phyxminiproblems-problemphyxmini0086-lengthinmeters}
  \lean{PhyXMiniProblems.ProblemPhyXMini0086.lengthInMeters}
  \uses{def:physics:phyx-mini-0086:phyxminiproblems-problemphyxmini0086-dimlength}
  The scalar meter readout of a dimensionful length
\end{definition}
\begin{proof}
  This is the declared model definition of length In Meters.
\end{proof}
\begin{definition}[time In Seconds]
  \label{def:physics:phyx-mini-0086:phyxminiproblems-problemphyxmini0086-timeinseconds}
  \lean{PhyXMiniProblems.ProblemPhyXMini0086.timeInSeconds}
  \uses{def:physics:phyx-mini-0086:phyxminiproblems-problemphyxmini0086-dimtime}
  The scalar second readout of a dimensionful time interval
\end{definition}
\begin{proof}
  This is the declared model definition of time In Seconds.
\end{proof}
\begin{definition}[time In Nanoseconds]
  \label{def:physics:phyx-mini-0086:phyxminiproblems-problemphyxmini0086-timeinnanoseconds}
  \lean{PhyXMiniProblems.ProblemPhyXMini0086.timeInNanoseconds}
  \uses{def:physics:phyx-mini-0086:phyxminiproblems-problemphyxmini0086-dimtime, def:physics:phyx-mini-0086:phyxminiproblems-problemphyxmini0086-nanosecondunitchoices}
  The scalar nanosecond readout of a dimensionful time interval
\end{definition}
\begin{proof}
  This is the declared model definition of time In Nanoseconds.
\end{proof}
\begin{definition}[speed In Meters Per Second]
  \label{def:physics:phyx-mini-0086:phyxminiproblems-problemphyxmini0086-speedinmeterspersecond}
  \lean{PhyXMiniProblems.ProblemPhyXMini0086.speedInMetersPerSecond}
  \uses{def:physics:phyx-mini-0086:phyxminiproblems-problemphyxmini0086-dimspeedreal}
  The scalar SI readout of a dimensionful speed, in meters per second
\end{definition}
\begin{proof}
  This is the declared model definition of speed In Meters Per Second.
\end{proof}
\begin{definition}[travel Time Difference Seconds]
  \label{def:physics:phyx-mini-0086:phyxminiproblems-problemphyxmini0086-traveltimedifferenceseconds}
  \lean{PhyXMiniProblems.ProblemPhyXMini0086.travelTimeDifferenceSeconds}
  \uses{def:physics:phyx-mini-0086:phyxminiproblems-problemphyxmini0086-opticalfiberpulsesetup, def:physics:phyx-mini-0086:phyxminiproblems-problemphyxmini0086-timeinseconds}
  The travel-time excess of the reflected route over the axial route, in seconds
\end{definition}
\begin{proof}
  This is the declared model definition of travel Time Difference Seconds.
\end{proof}
\begin{definition}[travel Time Difference Nanoseconds]
  \label{def:physics:phyx-mini-0086:phyxminiproblems-problemphyxmini0086-traveltimedifferencenanoseconds}
  \lean{PhyXMiniProblems.ProblemPhyXMini0086.travelTimeDifferenceNanoseconds}
  \uses{def:physics:phyx-mini-0086:phyxminiproblems-problemphyxmini0086-opticalfiberpulsesetup, def:physics:phyx-mini-0086:phyxminiproblems-problemphyxmini0086-timeinnanoseconds}
  The same travel-time excess, read in nanoseconds
\end{definition}
\begin{proof}
  This is the declared model definition of travel Time Difference Nanoseconds.
\end{proof}
\begin{definition}[Matches Optical Fiber Figure]
  \label{def:physics:phyx-mini-0086:phyxminiproblems-problemphyxmini0086-matchesopticalfiberfigure}
  \lean{PhyXMiniProblems.ProblemPhyXMini0086.MatchesOpticalFiberFigure}
  \uses{def:physics:phyx-mini-0086:phyxminiproblems-problemphyxmini0086-opticalfiberpulsesetup, def:physics:phyx-mini-0086:phyxminiproblems-problemphyxmini0086-lengthinmeters}
  Numerical and diagrammatic data supplied by the problem: n₁ = 1.58, n₂ = 1.53, L = 300 m, and the direct ray follows the central axis
\end{definition}
\begin{proof}
  This is the declared model definition of Matches Optical Fiber Figure.
\end{proof}
\begin{definition}[Satisfies Optical Fiber Laws]
  \label{def:physics:phyx-mini-0086:phyxminiproblems-problemphyxmini0086-satisfiesopticalfiberlaws}
  \lean{PhyXMiniProblems.ProblemPhyXMini0086.SatisfiesOpticalFiberLaws}
  \uses{def:physics:phyx-mini-0086:phyxminiproblems-problemphyxmini0086-opticalfiberregion, def:physics:phyx-mini-0086:phyxminiproblems-problemphyxmini0086-fiberroute, def:physics:phyx-mini-0086:phyxminiproblems-problemphyxmini0086-opticalfiberpulsesetup, def:physics:phyx-mini-0086:phyxminiproblems-problemphyxmini0086-lengthinmeters, def:physics:phyx-mini-0086:phyxminiproblems-problemphyxmini0086-timeinseconds, def:physics:phyx-mini-0086:phyxminiproblems-problemphyxmini0086-speedinmeterspersecond}
  The governing geometrical-optics and constant-speed laws for the two rays. The critical Snell relation is stated at a tangential transmitted angle. Unfolding the repeated specular reflections produces a straight segment whose axial projection is the fiber length, giving path * cos(angle-to-axis) = L. These laws do not state the requested time difference or any answer choice
\end{definition}
\begin{proof}
  This is the declared model definition of Satisfies Optical Fiber Laws.
\end{proof}
\begin{definition}[Answer Choice]
  \label{def:physics:phyx-mini-0086:phyxminiproblems-problemphyxmini0086-answerchoice}
  \lean{PhyXMiniProblems.ProblemPhyXMini0086.AnswerChoice}
  Labels of the four displayed multiple-choice answers
\end{definition}
\begin{proof}
  This is the declared model definition of Answer Choice.
\end{proof}
\begin{definition}[answer Choice Nanoseconds]
  \label{def:physics:phyx-mini-0086:phyxminiproblems-problemphyxmini0086-answerchoicenanoseconds}
  \lean{PhyXMiniProblems.ProblemPhyXMini0086.answerChoiceNanoseconds}
  \uses{def:physics:phyx-mini-0086:phyxminiproblems-problemphyxmini0086-answerchoice}
  The nanosecond value printed next to each answer choice
\end{definition}
\begin{proof}
  This is the declared model definition of answer Choice Nanoseconds.
\end{proof}
\begin{definition}[Is Closest Answer Choice]
  \label{def:physics:phyx-mini-0086:phyxminiproblems-problemphyxmini0086-isclosestanswerchoice}
  \lean{PhyXMiniProblems.ProblemPhyXMini0086.IsClosestAnswerChoice}
  \uses{def:physics:phyx-mini-0086:phyxminiproblems-problemphyxmini0086-answerchoice, def:physics:phyx-mini-0086:phyxminiproblems-problemphyxmini0086-answerchoicenanoseconds}
  A choice is uniquely closest to an actual nanosecond readout
\end{definition}
\begin{proof}
  This is the declared model definition of Is Closest Answer Choice.
\end{proof}
\begin{lemma}[critical reflection delay formula]
  \label{lem:physics:phyx-mini-0086:phyxminiproblems-problemphyxmini0086-critical-reflection-delay-formula}
  \lean{PhyXMiniProblems.ProblemPhyXMini0086.critical_reflection_delay_formula}
  \uses{def:physics:phyx-mini-0086:phyxminiproblems-problemphyxmini0086-opticalfiberpulsesetup, def:physics:phyx-mini-0086:phyxminiproblems-problemphyxmini0086-lengthinmeters, def:physics:phyx-mini-0086:phyxminiproblems-problemphyxmini0086-speedinmeterspersecond, def:physics:phyx-mini-0086:phyxminiproblems-problemphyxmini0086-traveltimedifferenceseconds, def:physics:phyx-mini-0086:phyxminiproblems-problemphyxmini0086-matchesopticalfiberfigure, def:physics:phyx-mini-0086:phyxminiproblems-problemphyxmini0086-satisfiesopticalfiberlaws}
  The critical-angle and projection laws give the exact excess travel time. This is an intermediate physical conclusion, not a premise of the answer
\end{lemma}
\begin{proof}
  The conclusion follows from the governing relations and figure readouts named in the statement of critical reflection delay formula.
\end{proof}
% --- Archon declaration topology end ---
% --- Archon physics formalization source end ---
